Un taquin est constitué de $15$ pièces numérotées et d'une case vide,
disposées dans une grille $4 \times 4$.

Une opération consiste à déplacer dans la case vide une pièce voisine
horizontalement ou verticalement.

On part de la configuration ordonnée
\[
\begin{array}{cccc}
1 & 2 & 3 & 4 \\
5 & 6 & 7 & 8 \\
9 & 10 & 11 & 12 \\
13 & 14 & 15 & \square
\end{array}
\]
Peut-on atteindre la configuration
\[
\begin{array}{cccc}
1 & 2 & 3 & 4 \\
5 & 6 & 7 & 8 \\
9 & 10 & 11 & 12 \\
13 & 15 & 14 & \square
\end{array}
\ ?
\]
\emph{Indications :}
\begin{itemize}
  \item chaque déplacement échange la case vide et une pièce ;
  \item colorier la grille comme un échiquier ;
  \item si la case vide revient à sa position initiale, le nombre de
  déplacements est pair ;
  \item comparer la parité de la permutation des $16$ objets.
\end{itemize}
